\scanpage{20}
\mspage{17}

Soit $x,y\in L$. $\exists$ chaîne $x_0,x_1,\ldots x_n$ avec $x_0=x$, $x_n=y$ et $x_i,x_{i+1}\in I_i^\circ$. \deleted{\unsure{Quitte}} Quitte à rapetisser $I_i$, \deleted{\unsure{Posons}} Quitte

Posons $y_i=\operatorname{or}(I_i)$, alors on a
\[
x_i\unsure{>|\geq}y_i\unsure{<|\leq}x_{i+1}
\unsure{>|\geq}y_{i+1}\unsure{<|\leq}x_{i+2}
\]
donc $L$ en tant qu'ens ord. est connexe : On peut connecter deux éléments quelconques par une chaîne d'inégalités. On va raisonner sur la longueur. Si la chaîne est de longueur $\leq1$, donc $x\leq y$ il suffit de \deleted{prouver} \inserted{noter} qu'il existe $z$ avec $y<z$ (provenant du fait que $y$ régulier). Supposons que c'est prouvé pour $n-1$, \deleted{donc $\exists$} prouvons le pour $n$. \deleted{\unsure{x}} On suppose donc $\exists z$ avec
\[
x\underset{n-1}{\sim}z,\qquad z\underset{1}{\sim}y.
\]
Par hypothèse de récurrence, $\exists x'$ tel que \deleted{$\unsure{x<}$} $x,z<x'$. Si $z\geq y$ alors $x'>x,y$ on gagne. Si $z<y$ \unsure{le seul qui a}

\mspage{19}

\rotatebox{90}{\unsure{EXTRA STRONG}}
\rotatebox{90}{\unsure{EXTRA STRONG}}
